\documentclass[11pt]{article}

%==============================================================
%  How src/model.py works — a code walkthrough
%==============================================================

\usepackage[a4paper,margin=1in]{geometry}
\usepackage{amsmath,amssymb,mathtools}
\usepackage{bm}
\usepackage{xcolor}
\usepackage{listings}
\usepackage{tcolorbox}
\tcbuselibrary{skins,breakable}
\usepackage{hyperref}
\usepackage{enumitem}
\usepackage{booktabs}
\usepackage{array}

\hypersetup{
  colorlinks=true,
  linkcolor=blue!50!black,
  urlcolor=blue!50!black,
}

%--- code styling -------------------------------------------------
\definecolor{codebg}{RGB}{248,248,248}
\definecolor{codekw}{RGB}{20,60,160}
\definecolor{codestr}{RGB}{40,120,40}
\definecolor{codecom}{RGB}{120,120,120}

\lstdefinestyle{py}{
  language=Python,
  basicstyle=\ttfamily\footnotesize,
  keywordstyle=\color{codekw}\bfseries,
  stringstyle=\color{codestr},
  commentstyle=\color{codecom}\itshape,
  backgroundcolor=\color{codebg},
  frame=single,
  framerule=0pt,
  rulecolor=\color{codebg},
  numbers=none,
  showstringspaces=false,
  breaklines=true,
  columns=fullflexible,
  keepspaces=true,
  xleftmargin=4pt,
  xrightmargin=4pt,
}
\lstset{style=py}

\newtcolorbox{ideabox}[1]{
  enhanced, breakable,
  colback=blue!2!white, colframe=blue!40!black,
  fonttitle=\bfseries,
  title={#1},
  boxrule=0.4pt, arc=2pt,
  left=4pt, right=4pt, top=4pt, bottom=4pt
}

\newcommand{\R}{\mathbb{R}}
\newcommand{\xv}{\mathbf{x}}
\newcommand{\uv}{\mathbf{u}}
\newcommand{\yv}{\mathbf{y}}
\newcommand{\Mv}{\mathbf{M}}
\newcommand{\Bv}{\mathbf{B}}
\newcommand{\Av}{\mathbf{A}}
\newcommand{\Fv}{\mathbf{F}}
\newcommand{\dd}{\mathrm{d}}
\newcommand{\code}[1]{\texttt{\small #1}}

\title{\textbf{How \texttt{src/model.py} works}\\
\large A walkthrough of the springs-and-sticks PyTorch SDE module}
\date{}
\author{}

\begin{document}
\maketitle

\noindent\code{src/model.py} defines two \code{nn.Module} classes that together implement the springs-and-sticks (SS) SDE in a form \code{torchsde.sdeint} can integrate:
\begin{itemize}[leftmargin=*,itemsep=2pt]
  \item \code{GroupGS3DE} --- a divide-and-conquer wrapper that splits the input features into groups and instantiates one core model per group.
  \item \code{GS3DE} --- the actual springs-and-sticks SDE on a single (sub-)domain. All of the interesting physics lives here.
\end{itemize}

\tableofcontents
\bigskip\hrule\bigskip

%==============================================================
\section{Architecture: \texttt{GroupGS3DE} wraps \texttt{GS3DE}}
%==============================================================

Building one symbolic Lagrangian over a full $d$-dimensional grid scales as $O(N_s^d)$ degrees of freedom --- exponential in the input dimension. \code{GroupGS3DE} mitigates this by partitioning the $d$ input features into smaller groups (of size $\le$ \code{max\_features}) and giving each group its own \code{GS3DE} model.

\begin{enumerate}[leftmargin=*,itemsep=2pt]
  \item \code{split\_features(num\_features, strategy)} returns a list of index tensors, either sequential (\code{torch.split(arange(d), max\_features)}) or random (\code{torch.split(torch.randperm(d), \dots)}).
  \item For each index group, a separate \code{GS3DE} is constructed with the corresponding slice of the boundaries. Construction can run in a \code{ThreadPoolExecutor} when \code{parallel=True}.
  \item \code{self.models} is an \code{nn.ModuleList}; \code{self.state\_sizes} tracks each sub-model's state size; \code{self.cstate\_sizes} is the cumulative sum, used to slice a joint \code{theta} into per-model chunks.
\end{enumerate}

All glue methods --- \code{f}, \code{g}, \code{cost}, \code{loss}, \code{y\_prediction}, \code{num\_y\_prediction} --- iterate over \code{self.models}, call the same method on each, and either concatenate (\code{f, g}), sum (\code{cost}), or average (\code{y\_prediction}) the results. \code{get\_theta\_group(theta, i)} returns the slice of \code{theta} that belongs to sub-model $i$:
\begin{lstlisting}
def get_theta_group(self, theta, group, predict=False):
    if predict:
        return theta[self.cstate_sizes[group]:self.cstate_sizes[group+1]]
    return theta[:, self.cstate_sizes[group]:self.cstate_sizes[group+1]]
\end{lstlisting}

The rest of this document is about \code{GS3DE}; \code{GroupGS3DE} is a thin orchestrator.

%==============================================================
\section{\texttt{GS3DE.\_\_init\_\_}: physical parameters and one-time setup}
%==============================================================

The constructor stores the physical parameters (\code{k}, \code{M}, \code{friction}, \code{temp}, \code{kb}, \code{k2}) and computes a few derived quantities used everywhere else:

\begin{lstlisting}
self.ell       = (self.u_max - self.u_min) / self.n_sticks
self.M         = M / torch.prod(self.n_sticks)
self.state_size = torch.prod(self.n_sticks + 1) * 2 * n_labels
self.eta_cte   = float(np.sqrt(2 * self.friction * temp * kb / self.M))
\end{lstlisting}

\noindent\begin{tabular}{@{}lp{0.74\linewidth}@{}}
\toprule
\code{ell}        & cell width per axis. \\
\code{M}          & total mass distributed evenly across sticks. \\
\code{state\_size} & $2$ phase-space coords (position + velocity) $\times$ $\prod_b(N_s^b+1)$ grid vertices $\times$ \code{n\_labels} output channels. This is the dimension \code{torchsde} integrates. \\
\code{eta\_cte}   & the noise amplitude $\sigma = \sqrt{2\gamma T k_B / M}$ from the fluctuation--dissipation relation. \\
\bottomrule
\end{tabular}

\medskip
Then it performs the \emph{symbolic} setup, once and for all, in fixed order:

\begin{enumerate}[leftmargin=*,itemsep=0pt]
  \item \code{\_init\_symbols} --- create sympy dynamic symbols on the grid.
  \item \code{\_init\_kinetic\_energy} --- build $K_\text{tr}, K_\text{rot}$.
  \item \code{\_init\_elastic\_energy} --- build $U_\text{el}$ (only nonzero if \code{k2 != 0}).
  \item \code{\_init\_lagrangian} --- assemble $L_\text{nonpot}$, derive the mass matrix and non-potential forcing, and call \code{\_update\_total\_lagrangian} once.
\end{enumerate}

\begin{ideabox}{The split: static skeleton vs.\ data-dependent piece}
The data-dependent potential $U(\xv)$ is \textbf{not} built in \code{\_\_init\_\_}: it is supplied later via \code{update\_data(u\_i, y\_i)}. Everything that does not depend on the data --- the symbolic positions, the kinetic energies, the mass matrix --- is computed exactly once. This is the whole point of the \code{\_init\_lagrangian} / \code{\_update\_total\_lagrangian} split.
\end{ideabox}

%==============================================================
\section{The symbolic skeleton}
%==============================================================

\subsection*{\texttt{\_init\_symbols}}

Builds a numpy array of \code{sympy.dynamicsymbols} with shape
\[
(\underbrace{N_s^0+1, \dots, N_s^{d-1}+1}_{\text{grid}},\, \underbrace{m}_{\text{output channels}}).
\]
So \code{x\_symbols[i, j, ..., p]} is the height $x_{(i,j,\dots)}^p$ of the stick at that grid vertex in output channel $p$. The arrays \code{dx\_symbols} and \code{ddx\_symbols} hold the first and second time-derivatives of each symbol. Total number of symbolic variables: \code{self.N = prod(shape)}.

\subsection*{\texttt{\_init\_kinetic\_energy}}

For each axis $b$ it takes the ``front'' slice (\code{slice(1, None)}) and ``back'' slice (\code{slice(None, -1)}) along that axis to extract neighboring stick velocity pairs $(\dot x_{\vec i + \hat e^b}, \dot x_{\vec i})$:

\begin{lstlisting}
for i in range(len(shape) - 1):                # over input dims
    sf = [slice(None)]*len(shape); sf[i] = slice(1, None)
    sb = [slice(None)]*len(shape); sb[i] = slice(None, -1)
    diff = (dx_symbols[tuple(sf)] + dx_symbols[tuple(sb)])**2
    ktr += (M/8) * Add(*diff.ravel())          # translational
    diff = (dx_symbols[tuple(sf)] - dx_symbols[tuple(sb)])**2
    krot += (M/24) * Add(*diff.ravel())        # rotational
\end{lstlisting}

The result is sympy expressions:
\[
K_\text{tr} = \frac{M}{8}\sum_{b,\vec i}(\dot x_{\vec i+\hat e^b} + \dot x_{\vec i})^2,
\qquad
K_\text{rot} = \frac{M}{24}\sum_{b,\vec i}(\dot x_{\vec i+\hat e^b} - \dot x_{\vec i})^2.
\]

\subsection*{\texttt{\_init\_elastic\_energy}}

Same neighbor pattern but on positions, with an offset by \code{self.ell[i]}:
\[
U_\text{el} = \frac{k_2}{2}\sum_{b,\vec i}\bigl(x_{\vec i+\hat e^b} - x_{\vec i} - \ell^b\bigr)^2.
\]
Skipped (left as $0$) when \code{k2 == 0}, which is the default.

\subsection*{\texttt{\_init\_lagrangian}}

Combines the non-potential pieces and runs sympy's Euler--Lagrange machinery:

\begin{lstlisting}
self.L_nonpot = simplify(Add(self.ktr, self.krot, -self.uelastic_symbols))
LM = LagrangesMethod(self.L_nonpot, self.x_symbols.flatten())
LM.form_lagranges_equations()
self.mass_matrix     = simplify(LM.mass_matrix)
self.inv_mass_matrix = simplify(LM.mass_matrix.inv())   # falls back to .pinv()
self.forcing_nonpot  = simplify(LM.forcing)
self._update_total_lagrangian()                         # see next section
self.ypred = lambda u: lambdify([*x_symbols.flatten()],
                                 self.y_prediction(u))
\end{lstlisting}

\code{self.mass\_matrix} and \code{self.inv\_mass\_matrix} live for the lifetime of the model: they don't depend on data.

%==============================================================
\section{Bringing in data: \texttt{update\_data}}
%==============================================================

\code{update\_data(new\_u\_i, new\_y\_i)} does exactly two things:

\subsection*{\texttt{\_compute\_potential\_energy(u\_i, y\_i)}}
\begin{enumerate}[leftmargin=*,itemsep=0pt]
  \item \code{find\_box(u\_i)} returns the integer multi-index $\vec i$ of the cell that contains $\uv_i$.
  \item \code{y\_prediction(u\_i, i\_idx)} returns the symbolic vector $\hat\yv(\uv_i)$ as a multilinear combination of the surrounding stick heights (see Section~\ref{sec:pwlin}).
  \item Builds the symbolic potential
  \[
  U(\xv) \;=\; \frac{k}{2}\sum_{j,p}\bigl(\hat y^p(\uv_j) - y_j^p\bigr)^2.
  \]
\end{enumerate}

\subsection*{\texttt{\_update\_total\_lagrangian}}
\begin{lstlisting}
self.lagrangian    = Add(self.L_nonpot, -self.U)
self.forcing_pot   = Matrix([-diff(self.U, q) for q in x_symbols.flatten()])
self.forcing_vector = Add(self.forcing_nonpot, self.forcing_pot)
self.evol_dynamics = simplify(self.inv_mass_matrix * self.forcing_vector)

self.lambdified_dyn = lambdify(
    [*x_symbols.flatten(), *dx_symbols.flatten()],
    self.evol_dynamics - self.friction * dx_symbols.reshape(-1, 1))
self.ue = lambdify([*x_symbols.flatten()], self.U)
\end{lstlisting}

After this call:
\begin{itemize}[leftmargin=*,itemsep=0pt]
  \item \code{lambdified\_dyn(*q, *dq)} returns $\ddot\xv = \Mv^{-1}\Fv(\xv,\dot\xv) - \gamma\dot\xv$ as fast numerical code (compiled by sympy's \code{lambdify}).
  \item \code{ue(*q)} returns $U(\xv)$ numerically.
\end{itemize}

These two compiled functions are what \code{f}, \code{cost}, \code{dw} actually call at integration time --- no sympy in the hot loop.

%==============================================================
\section{Piecewise-linear prediction}\label{sec:pwlin}
%==============================================================

Three helpers turn an input $\uv$ into a prediction:

\begin{lstlisting}
def find_box(self, u):
    i = ((u - self.u_min) / self.ell).to(torch.int)
    return torch.clip(i, 0, self.n_sticks - 1)

def lamd(self, i, u, flip_ind=None, div_by_ell=True):
    val = u - i*self.ell - self.u_min
    if div_by_ell: val /= self.ell
    return val                              # the lambda_i^b(u)
\end{lstlisting}

\code{find\_box} maps $\uv$ to the integer cell index $\vec i$. \code{lamd} returns the in-cell coordinate $\lambda_{\vec i}^b(\uv) \in [0,1]$.

\code{\_y\_prediction(u, i=None)} assembles the multilinear interpolant:
\begin{lstlisting}
u = torch.clamp(u, u_min, u_max)
i = find_box(u) if i is None else i
ld = lamd(i, u)
offsets = torch.tensor(list(product([0,1], repeat=d)))
grid_points = i.unsqueeze(1) + offsets.unsqueeze(0)
weights = torch.prod((1-ld)*(offsets==0) + ld*(offsets!=0), axis=2)
assert torch.all(weights.sum(axis=1) - 1 < 1e-5)
symbols = self.x_symbols[index_tuple]
pred = sum(weights * symbols, axis=corners)
\end{lstlisting}

The assertion is the sanity check that $\sum_{\vec j} w_{\vec j} = 1$ inside every cell. The result \code{pred} is a vector of \emph{symbolic} expressions in the stick positions --- ready to be squared into a potential.

\code{num\_y\_prediction(u, theta)} is the purely numerical sibling: it evaluates the cached \code{self.ypred(u)} (which is \code{lambdify}d from the symbolic prediction) on the positions stored in \code{theta}.

%==============================================================
\section{The SDE interface: \texttt{f} and \texttt{g}}
%==============================================================

These are the two methods \code{torchsde.sdeint} calls at every Euler step. The argument \code{theta} has shape \code{(batch, state\_size)} with the first half holding positions $\xv$ and the second half velocities $\dot\xv$.

\subsection*{\texttt{f(t, theta)} --- the drift}
\begin{lstlisting}
def f(self, t, theta):
    q  = theta[:, :self.N].t()                  # positions
    dq = theta[:, self.N:].t()                  # velocities
    ddqdt = torch.tensor(self.lambdified_dyn(*q, *dq))
    return torch.vstack([dq, ddqdt]).t()
\end{lstlisting}
This computes
\[
\mathbf{f}(\boldsymbol\theta) \;=\; \begin{pmatrix} \dot\xv \\ \Mv^{-1}\Fv(\xv,\dot\xv) - \gamma\dot\xv \end{pmatrix}.
\]
The branch on \code{ddqdt\_shape} handles the degenerate one-dimensional case.

\subsection*{\texttt{g(t, theta)} --- the diffusion}
\begin{lstlisting}
def g(self, t, theta):
    return torch.cat([
        torch.zeros_like(theta[:, :self.N]),
        self.eta_cte * torch.ones_like(theta[:, self.N:])
    ], dim=1)
\end{lstlisting}
Position rows are zero (positions inherit noise only \emph{indirectly}, through integrating noisy velocities); velocity rows are \code{eta\_cte} $=\sigma$. Combined with \code{noise\_type = "diagonal"} and \code{sde\_type = "ito"}, this tells \code{torchsde} the SDE is
\[
\dd\boldsymbol\theta \;=\; \mathbf{f}(\boldsymbol\theta)\,\dd t \;+\; \Bv\,\dd W_t,
\qquad \Bv = \mathrm{diag}(0,\dots,0,\sigma,\dots,\sigma).
\]

%==============================================================
\section{Loss, energy, work}
%==============================================================

\begin{tabular}{@{}lp{0.78\linewidth}@{}}
\toprule
\code{cost(theta)}  & Evaluates \code{self.ue(*q)} --- the data potential $U(\xv)$ on each batch element. Since $U=(k/2)\sum(\hat y - y)^2$, this is the (unnormalized) squared-error loss times $k/2$. \\
\code{loss(theta, u\_i, y\_i)} & Returns $2U/(N\cdot k) =$ MSE per data point. \\
\code{dcost(theta)} & The total time-derivative $\dd U/\dd t$, computed by sympy via \code{self.U.diff()} and then \code{lambdify}d. \\
\code{dw(theta)} & The instantaneous power $\dd W/\dd t = M\langle\ddot\xv, \dot\xv\rangle$ --- used by the Jarzynski free-energy estimator in \code{src/entropy.py}. \\
\bottomrule
\end{tabular}

\begin{lstlisting}
def dw(self, theta):
    q  = theta[:, :self.N].t()
    dq = theta[:, self.N:].t()
    ddqdt = torch.tensor(self.lambdified_dyn(*q, *dq)).squeeze()
    return torch.sum(ddqdt * dq, dim=0) * self.M
\end{lstlisting}

%==============================================================
\section{Lifecycle in a training loop}
%==============================================================

The usage pattern (as in \code{src/profiling.py}) is:

\begin{lstlisting}
sde = GroupGS3DE(max_features, n_sticks, boundaries, n_labels, ...)
theta0 = None
for u_i, y_i in loader_train:
    sde.update_data(u_i, y_i)                # rebuild U, ddqdt, ue
    if theta0 is None:
        theta0 = torch.rand((num_runs, sde.state_size))
    thetas = torchsde.sdeint(sde, theta0, ts, method='euler')
    # sde.f and sde.g are called at every Euler step
    theta0 = thetas[-1, :]                   # warm-start next batch
\end{lstlisting}

Only data-dependent pieces are rebuilt between mini-batches. The kinetic energy, mass matrix, and elastic energy are computed exactly once in \code{\_\_init\_\_}; the data potential and the compiled drift function are rebuilt on every \code{update\_data}; the SDE solver only ever calls the cheap numerical \code{f, g}.

%==============================================================
\section{How are the sticks actually moved forward in time?}
%==============================================================

There is a tiny time increment $\Delta t$ and the integration scheme is \textbf{Euler--Maruyama}, but \code{model.py} does not run the time-stepping itself --- it just hands \code{f} and \code{g} to \code{torchsde.sdeint}, which loops in C.

\medskip
\noindent The state vector is
\[
\boldsymbol\theta(t) \;=\; \begin{pmatrix} \xv(t) \\ \dot\xv(t) \end{pmatrix} \in \R^{2N}.
\]
One Euler--Maruyama step is
\[
\boldsymbol\theta(t+\Delta t) \;=\; \boldsymbol\theta(t) \;+\; \mathbf{f}(t,\boldsymbol\theta)\,\Delta t \;+\; \mathbf{g}(t,\boldsymbol\theta)\,\sqrt{\Delta t}\;\bm\eta, \qquad \bm\eta \sim \mathcal{N}(0,\mathbb{I}).
\]
Concretely, what \code{torchsde.sdeint(sde, theta0, ts, method='euler')} does is essentially:

\begin{lstlisting}
theta = theta0
for k in range(len(ts) - 1):
    dt    = ts[k+1] - ts[k]
    eta   = torch.randn_like(theta)
    drift = sde.f(ts[k], theta)              # (dx/dt, d2x/dt2)
    noise = sde.g(ts[k], theta)              # (0,...,0, sigma,...,sigma)
    theta = theta + drift*dt + noise*eta*sqrt(dt)
\end{lstlisting}

So \code{sde.f} returns \emph{velocities and accelerations} (how fast positions and velocities are currently changing), and the solver multiplies that by \code{dt} to push the state forward one tick. The accelerations themselves are computed inside \code{sde.f} by calling \code{self.lambdified\_dyn(*q, *dq)}, which is a numerically compiled function that evaluates
\[
\ddot\xv \;=\; \Mv^{-1}\Fv(\xv,\dot\xv) - \gamma\dot\xv
\]
at the current state.

\begin{ideabox}{Where does $\Delta t$ come from?}
$\Delta t$ comes entirely from the user's choice of \code{ts = torch.linspace(t\_start, t\_end, t\_size)}; the solver uses the gaps between successive entries of \code{ts}. In \code{profiling.py} you have \code{ts = torch.linspace(0, 1, 2)}, which means one Euler step of size $\Delta t = 1$ (good for profiling, bad for accuracy). For real runs \code{t\_size} should be in the hundreds or thousands so $\Delta t$ is small.
\end{ideabox}

\noindent\textbf{Take-away.}\quad The sticks ``move'' inside \code{torchsde.sdeint}, not in any \code{model.py} function. \code{model.py} only specifies the \emph{rules} of motion (\code{f} for the drift, \code{g} for the noise); the solver applies those rules step by step.

%==============================================================
\section{Function-wise flow at a glance}
%==============================================================

\begin{lstlisting}[basicstyle=\ttfamily\scriptsize]
profiling.run_main()
|
+-- GroupGS3DE.__init__()                            [ONCE, ~seconds-minutes]
|   +-- GS3DE.__init__()         per feature group
|       +-- _init_symbols()                    <- sympy variables
|       +-- _init_kinetic_energy()             <- K_tr, K_rot
|       +-- _init_elastic_energy()             <- U_el (often 0)
|       +-- _init_lagrangian()
|           +-- LagrangesMethod(...)           <- mass_matrix, inv_mass_matrix
|           +-- _update_total_lagrangian()     <- no-op: U is still 0
|
+-- for each epoch:
    +-- for each mini-batch (u_i, y_i):
        |
        +-- sde.update_data(u_i, y_i)                [PER BATCH]
        |   +-- _compute_potential_energy(u_i, y_i)
        |   |   +-- find_box(u_i)                    <- which cell?
        |   |   +-- y_prediction(u_i, i_idx)         <- symbolic y_hat
        |   |       +-- _y_prediction()
        |   |           +-- lamd(i, u)               <- lambda in [0,1]
        |   |           +-- (multilinear corner weights)
        |   +-- _update_total_lagrangian()
        |       +-- lagrangian   = L_nonpot - U
        |       +-- evol_dynamics = M^-1 * forcing
        |       +-- lambdified_dyn = lambdify(...)   <- FAST numeric ddx(x,dx)
        |
        +-- sde.loss(theta0, u_i, y_i)               <- report MSE
        |
        +-- torchsde.sdeint(sde, theta0, ts, 'euler') [PER EULER STEP x t_size]
        |   +-- for k in range(t_size - 1):
        |       +-- sde.f(t, theta)
        |       |   +-- lambdified_dyn(*q, *dq)      <- compute ddx
        |       +-- sde.g(t, theta)                  <- sigma on velocity rows
        |       +-- theta += f*dt + g*eta*sqrt(dt)   <- Euler-Maruyama
        |
        +-- theta0 = thetas[-1]                      <- warm-start next batch
\end{lstlisting}

\subsection*{How to read this tree}

\begin{itemize}[leftmargin=*,itemsep=2pt]
  \item \textbf{Top block (\code{\_\_init\_\_})} runs once and builds the symbolic skeleton that does \emph{not} depend on data: mass matrix, inverse mass matrix, kinetic energies. This is the slow sympy work; it is amortized over the whole training run.
  \item \textbf{\code{update\_data} (per batch)} rebuilds only the data-dependent piece $U(\xv)$ and recompiles the numerical drift function \code{lambdified\_dyn}. The mass matrix is \emph{not} rebuilt --- it never changes.
  \item \textbf{\code{sdeint} (per Euler step inside a batch)} is the hot loop. It only ever calls the cheap numerical functions \code{f} and \code{g} --- no sympy in this loop. Each call advances the state by one $\Delta t$.
\end{itemize}

%==============================================================
\section{Summary table}
%==============================================================

\noindent\begin{tabular}{p{0.34\linewidth} p{0.42\linewidth} p{0.18\linewidth}}
\toprule
\textbf{Method} & \textbf{Role} & \textbf{Recomputed} \\
\midrule
\code{\_init\_symbols}        & sympy variables on the grid                              & once \\
\code{\_init\_kinetic\_energy}  & $K_\text{tr}, K_\text{rot}$ symbolically                 & once \\
\code{\_init\_elastic\_energy}  & $U_\text{el}$ symbolically (if \code{k2}$\ne 0$)         & once \\
\code{\_init\_lagrangian}       & $L_\text{nonpot}$, mass matrix, non-pot.\ forcing        & once \\
\code{\_compute\_potential\_energy} & build $U(\xv)$ from new data                          & per \code{update\_data} \\
\code{\_update\_total\_lagrangian}  & rederive $\ddot\xv$ symbolically, lambdify it          & per \code{update\_data} \\
\code{find\_box, lamd, \_y\_prediction} & piecewise-linear interpolation kernel             & per call \\
\code{f(t, theta)}            & SDE drift for \code{torchsde}                          & per Euler step \\
\code{g(t, theta)}            & SDE noise matrix                                       & per Euler step \\
\code{cost, loss, dcost, dw}  & diagnostics (energy, MSE, power)                       & on demand \\
\bottomrule
\end{tabular}

\end{document}
